\subsection{State estimation for a moving dock}
\label{sec:lit-estimation}

When the dock moves, the guidance layer of the control system needs a model of the dock's motion in addition to its pose. The target-tracking literature supplies a standard kinematic hierarchy for this \cite{barshalom2001estimation}. A constant-position model treats the target as static and absorbs all dock relative motion as process noise on the position. A nearly-constant-velocity model augments an additional velocity state driven by white acceleration noise, and manoevre models add further derivatives or switching. The acceleration noise density is the main tuning parameter as it sets the filter's bandwidth, which sets the filter's bandwidth as a trade-off between responsiveness to real target motion and amplification of measurement noise. In regards to this thesis, two properties of this hierarchy matter. First, the constant-position model is a different estimator since it does not have a velocity state and cannot extrapolate, which limits its tracking against a moving target, but immune to velocity estimation corruption. Second, evert member of the hierarchy inherits the Kalman assumption that measurement errors are white noises and independent of the state. The filter offers no guard against this false assumption and no process-noise setting can retune it away \cite{barshalom2001estimation}. The two filter properties are compared in Part~A and its failure analysis is a case study of the second property.

This kind of filter now appear routinely inside docking systems. For example, G\"unzel's resident-vehicle system fuses ArUco pose with acoustic positioning via an extended Kalman filter (eKf) \cite{gunzel2026docking}, and Ni's system fuses two visual pose changes by confidence weighting \cite{ni2025vision}. In both systems however, the dock is static, and the filter suppresses noise and outliers rather than estimating target motion. As a result, tracking a moving dock from a moving vehicle raises further problems. The marker measurement is reconstructed in the world frame through the vehicle's own estimated pose, so vehicle motion and localisation error leaks into the measured dock position. Therefore, this measurement error becomes correlated with the observer's state. This correlated measurement noise is a recognised failure more in target-tracking literature, with known flaws such as augmenting the filter with the observer's state or estimating in relative frame \cite{barshalom2001estimation}. In the underwater docking literature, however, this fail mode exposes an unexamined gap. Against a static dock, the leakage is harmless and absorbed as position noise, whereas a velocity state integrates it into a target velocity. Part~A demonstrates and analyses this mechanism, including the case where the corrupted velocity estimate feeds a control loop.
